\documentclass[11pt,a4paper]{article}

\usepackage[utf8]{inputenc}
\usepackage[T1]{fontenc}
\usepackage[spanish,es-tabla]{babel}
\usepackage[a4paper,margin=2.5cm]{geometry}
\usepackage{booktabs}
\usepackage{longtable}
\usepackage{array}
\usepackage{xcolor}
\usepackage{enumitem}
\usepackage{hyperref}
\usepackage{fancyhdr}
\usepackage{titlesec}

\definecolor{riesgoalto}{RGB}{178,34,34}
\definecolor{riesgomedio}{RGB}{184,134,11}
\definecolor{gris}{RGB}{90,90,90}

\hypersetup{
    colorlinks=true,
    linkcolor=blue!50!black,
    urlcolor=blue!50!black,
    pdftitle={Evaluacion tecnica de Mammouth.ai},
    pdfauthor={Escaneria - Migasa}
}

\titleformat{\section}{\normalfont\Large\bfseries}{\thesection}{1em}{}
\titleformat{\subsection}{\normalfont\large\bfseries}{\thesubsection}{1em}{}

\pagestyle{fancy}
\fancyhf{}
\fancyhead[L]{\small Evaluación técnica --- Mammouth.ai}
\fancyhead[R]{\small Proyecto Escaneria (Migasa)}
\fancyfoot[C]{\thepage}

\newcommand{\riesgo}[1]{\textcolor{riesgoalto}{\textbf{[#1]}}}
\newcommand{\advertencia}[1]{\textcolor{riesgomedio}{\textbf{[#1]}}}

\title{\textbf{Evaluación técnica de Mammouth.ai como capa de acceso\\y orquestación de modelos de IA} \\[0.3em]
\large Proyecto de extracción automática de facturas (Escaneria) --- Migasa}
\author{}
\date{11 de agosto de 2026}

\begin{document}

\maketitle
\thispagestyle{empty}

\begin{abstract}
\noindent
Este informe evalúa la viabilidad de usar \textbf{Mammouth.ai} como capa de acceso y
orquestación frente a distintos proveedores de modelos de IA, en sustitución o como
respaldo de la integración directa actual con la API de OpenAI en el pipeline de
extracción automática de datos de facturas. Se analiza el encaje con la arquitectura
existente, las capacidades de integración documentadas, los riesgos técnicos y
económicos, y dos escenarios de uso: sustitución del proveedor principal y uso como
contingencia ante una caída de OpenAI.
\end{abstract}

\tableofcontents

\section{Contexto y objetivo}

El proyecto Escaneria procesa facturas de proveedor de forma automática: extrae
texto de los PDF, y cuando el texto no es fiable o el documento es una imagen
escaneada, recurre a visión por computador sobre el propio PDF. Toda la
inteligencia de extracción se apoya hoy en la API de OpenAI, a través de dos
funciones centralizadas en \texttt{logic.py}:

\begin{itemize}
    \item \texttt{build\_client()} --- crea el cliente del SDK de OpenAI leyendo
    \texttt{OPENAI\_API\_KEY} y, opcionalmente, \texttt{OPENAI\_BASE\_URL} (pensado
    para poder apuntar a un endpoint alternativo compatible sin tocar el resto
    del código).
    \item \texttt{get\_model()} --- devuelve el modelo a usar, configurable vía
    \texttt{OPENAI\_MODEL} (por defecto \texttt{gpt-4.1}).
\end{itemize}

El motivo por el que se plantea evaluar Mammouth.ai es doble:

\begin{enumerate}
    \item Evitar depender de integraciones directas con cada proveedor de IA si en
    el futuro se quisiera usar más de un modelo.
    \item \textbf{Disponer de una vía alternativa si un día OpenAI deja de estar
    disponible}, para poder seguir procesando facturas con otro proveedor de IA.
\end{enumerate}

Este segundo punto es, tras el análisis, el que realmente encaja con el caso de
uso: se evalúa por tanto tanto el escenario de \emph{sustitución} como el de
\emph{contingencia} (sección~\ref{sec:contingencia}).

\section{Qué es Mammouth.ai}

Mammouth AI SAS es una empresa francesa (sede en París), fundada en 2024, con una
plantilla pequeña (2-10 empleados). Su producto principal es una
\textbf{suscripción de consumo/equipo} que da acceso, bajo un único login, a
más de 25 modelos de terceros (GPT, Claude, Gemini, Mistral, DeepSeek, GLM,
generadores de imagen y vídeo, etc.). El acceso programático mediante API es una
funcionalidad \emph{secundaria} de ese producto, no su núcleo de negocio.

\subsection{Capacidades de integración documentadas}

\begin{itemize}
    \item Endpoint: \texttt{https://api.mammouth.ai/v1/chat/completions}, compatible
    con el formato de la API de OpenAI (Chat Completions).
    \item Autenticación: cabecera \texttt{Authorization: Bearer <API\_KEY>}.
    \item Streaming de respuestas (\texttt{stream: true}) vía Server-Sent Events.
    \item Embeddings (\texttt{text-embedding-3-large/small}).
    \item Selección de modelo por nombre o por alias de conveniencia
    (\texttt{mammouth-recommended}, \texttt{mammouth}, \texttt{recommended}) que
    elige modelo automáticamente y puede cambiar sin previo aviso.
    \item Compañía europea, declara cumplimiento GDPR y no usar contenido de
    usuario para entrenar modelos; se financia solo por suscripción.
\end{itemize}

\subsection{Lo que la documentación pública \emph{no} confirma}

\begin{itemize}
    \item Soporte de \texttt{response\_format} con JSON Schema en modo
    \texttt{strict} (Structured Outputs).
    \item Soporte de imágenes como \emph{entrada} (visión/comprensión de imagen);
    solo documentan modelos de \emph{generación} de imagen.
    \item Límites de peticiones por minuto/hora para uso con clave de API.
    \item SLA de disponibilidad formal.
    \item Tarifa por token para uso de pago por consumo (\emph{pay-as-you-go}).
\end{itemize}

\section{Viabilidad técnica frente a la arquitectura actual}

El coste de \emph{conectar} Mammouth es prácticamente nulo: al existir ya
\texttt{OPENAI\_BASE\_URL} en \texttt{build\_client()}, probarlo es cambiar dos
variables en \texttt{.env}. El problema no es la conexión, sino que dos piezas
muy concretas del pipeline dependen de comportamiento \emph{no estándar} de la
API de OpenAI:

\begin{enumerate}
    \item \riesgo{Riesgo crítico} \textbf{Structured Outputs estricto.}
    \texttt{extract\_invoice\_with\_agent} (\texttt{logic.py}) y la extracción por
    visión (\texttt{imagenes.py}) usan
    \texttt{response\_format={"type": "json\_schema", "strict": true}}. Este modo
    estricto es precisamente lo que corrigió un bug histórico del proyecto (el
    CIF del proveedor terminando en la columna de TipoIVA por desplazamiento de
    campos). Si Mammouth no respeta ese modo con la misma garantía, se podría
    reintroducir ese bug sin que salte ningún error.
    \item \riesgo{Riesgo crítico} \textbf{Visión sobre imágenes de entrada.}
    \texttt{imagenes.py} manda páginas del PDF como \texttt{image\_url} en base64
    para leer el CIF en facturas escaneadas sin texto extraíble. No hay
    confirmación de que la capa de compatibilidad de Mammouth soporte
    comprensión de imagen de entrada, solo generación.
    \item \advertencia{Riesgo medio} \textbf{No determinismo de modelo.} Los
    alias de selección automática de modelo pueden cambiar de versión sin
    aviso, lo que dificulta auditar por qué una factura concreta se extrajo mal.
\end{enumerate}

\textbf{Recomendación de verificación:} antes de tomar cualquier decisión, replicar
exactamente la llamada de texto (con el \texttt{json\_schema strict}) y la llamada
de visión contra el endpoint de Mammouth, usando una factura con datos ya
conocidos (\emph{ground truth}), y comparar el resultado campo a campo.

\section{Ventajas}

\begin{itemize}
    \item Coste de integración casi nulo: la abstracción ya existente
    (\texttt{build\_client()}/\texttt{get\_model()}) permite probarlo sin tocar
    lógica de negocio.
    \item Compañía con sede en la Unión Europea, con compromiso GDPR declarado y
    modelo de negocio basado en suscripción (no en monetizar datos de usuario).
    \item Un único panel/factura si en algún momento se quisiera tener acceso a
    varios modelos (Claude, Gemini) sin gestionar N claves y N SDKs distintos.
\end{itemize}

\section{Desventajas y riesgos}

\subsection{Técnicos}
\begin{itemize}
    \item Compatibilidad no confirmada con \texttt{response\_format strict} (sección~3).
    \item Compatibilidad no confirmada con visión de entrada (sección~3).
    \item Selección de modelo no determinista mediante alias.
\end{itemize}

\subsection{Modelo de cuotas, pensado para chat interactivo, no para un backend 24/7}
\begin{itemize}
    \item Cuotas por \emph{sesión}, renovadas cada 3 horas.
    \item Al agotar la cuota del modelo elegido, Mammouth \textbf{degrada
    automáticamente a un modelo más ligero, sin lanzar ningún error} (por
    ejemplo, Opus~$\rightarrow$~Sonnet~$\rightarrow$~Haiku). Para un proceso
    desatendido que procesa lotes de facturas, esto significa degradación
    silenciosa de la calidad de extracción, no un límite que se pueda capturar
    y reintentar.
    \item Sin límites de peticiones por minuto/hora documentados para uso con
    clave de API.
    \item Créditos incluidos no acumulables entre meses.
\end{itemize}

\subsection{Económicos}
Ver análisis detallado en la sección~\ref{sec:coste}. En resumen: al volumen real
de Migasa, el crédito incluido en cualquiera de los planes se agota en los
primeros días del mes.

\subsection{Negocio y continuidad}
\begin{itemize}
    \item Empresa muy pequeña (2-10 empleados), fundada en 2024, sin SLA de
    disponibilidad publicado.
    \item Reseñas independientes reportan soporte deficiente y dificultades
    para cancelar o conseguir reembolsos.
    \item El acceso vía API es una funcionalidad secundaria de un producto
    orientado a consumidor final, no el foco de su negocio.
\end{itemize}

\subsection{Datos y cumplimiento}
\begin{itemize}
    \item Se añade un salto adicional de tratamiento de datos: CIF/NIF e
    importes de Migasa y sus proveedores pasarían por Mammouth antes de llegar
    al proveedor real del modelo (OpenAI, Anthropic, Google...).
    \item No hay un DPA (Data Processing Agreement) formal verificado para este
    caso de uso concreto, más allá de la política de privacidad genérica.
\end{itemize}

\section{Análisis de consumo de tokens y coste}
\label{sec:coste}

\subsection{Estimación por tipo de llamada}

La siguiente tabla se basa en los límites que ya existen en el propio código
(\texttt{MAX\_CHARS = 35000} en \texttt{logic.py}, \texttt{DPI\_IMAGENES = 200} y
\texttt{MAX\_PAGINAS = 3} en \texttt{imagenes.py}), no en cifras de marketing.

\begin{longtable}{@{}p{4cm}p{5.5cm}p{2.3cm}p{1.8cm}@{}}
\toprule
\textbf{Llamada} & \textbf{Origen del tamaño} & \textbf{Input (tokens)} & \textbf{Output} \\
\midrule
\endhead
\texttt{detectar\_numero\_factura} (por página) &
Prompt fijo corto + página de texto (hasta 15.000 caracteres) &
500 (típico) -- 3.750 (límite) & 5--10 \\
\addlinespace
\texttt{extract\_invoice\_with\_agent} (por factura) &
\texttt{DEFAULT\_PROMPT} ($\approx$2.100 tokens fijos) + texto de la factura
(hasta 35.000 caracteres) + JSON Schema de 19 campos ($\approx$500--700 tokens) &
3.450 (típico) -- 11.500 (límite) & 150--300 \\
\addlinespace
Fallback de visión (hasta 3 páginas a 200~DPI) &
3 imágenes tamaño A4 ($\approx$1.100 tokens/imagen) + prompt y schema &
6.000 -- 6.500 & 200--300 \\
\bottomrule
\end{longtable}

\subsection{Coste equivalente a tarifa de OpenAI directa}

Usando la tarifa pública de \texttt{gpt-4.1} (\$2/millón de tokens de entrada,
\$8/millón de salida, el modelo por defecto de \texttt{get\_model()}):

\begin{itemize}
    \item Factura simple con texto extraíble (1 página): \textbf{\$0,009 -- \$0,010} por factura.
    \item Factura larga/compleja (cerca del límite de 35.000 caracteres): \textbf{\$0,025 -- \$0,03} por factura.
    \item Factura escaneada con fallback de visión: \textbf{+\$0,015 -- \$0,02} adicionales sobre el intento de texto previo.
\end{itemize}

\subsection{Datos reales de producción}

El consumo real reportado es de \textbf{\$3/día} con una media de
\textbf{400 facturas/día}, es decir:

\begin{itemize}
    \item \textbf{\$0,0075 por factura} (coincide con la parte baja del rango estimado).
    \item \textbf{$\approx$\$90/mes} ($\approx$\$1.080/año).
    \item \textbf{$\approx$12.000 facturas/mes.}
\end{itemize}

\subsection{Comparativa con los planes de Mammouth}

\begin{table}[h]
\centering
\begin{tabular}{lccc}
\toprule
\textbf{Plan} & \textbf{Precio/mes} & \textbf{Crédito incluido} & \textbf{Cobertura del consumo real} \\
\midrule
Starter  & \$12 & \$2  & $\approx$16 horas de 720 \\
Standard & \$24 & \$4  & $\approx$32 horas \\
Expert   & \$72 & \$10 & $\approx$3,3 días de 30 \\
\bottomrule
\end{tabular}
\caption{Horas/días del mes que cubre el crédito incluido de cada plan al ritmo
real de consumo de Migasa (\$3/día).}
\end{table}

Incluso en el plan más caro para un solo usuario, el crédito incluido se agota en
los primeros 3--4 días del mes. El resto de días, cada mes, habría que recurrir a
pago adicional por consumo (tarifa no publicada) o asumir la degradación
automática de modelo.

\subsection{Coste total estimado si se usara como motor principal}

Suponiendo el escenario \emph{más favorable posible} para Mammouth --- que su
tarifa de pago por consumo fuera idéntica a la de OpenAI directo, sin ningún
margen de reventa ---:

\begin{align*}
\text{Plan Expert} &= \$72/\text{mes} \\
+\ \text{consumo real} &= \$90/\text{mes} \\
-\ \text{crédito ya incluido} &= -\$10/\text{mes} \\
\hline
\textbf{Total estimado} &\approx \textbf{\$152/mes}
\end{align*}

Esto supone un \textbf{$\approx$69\% más caro} que el coste actual (\$90/mes) por
el mismo resultado, en el caso optimista y sin degradación de modelo. Con
cualquier margen de reventa razonable, el coste real probablemente estaría entre
\$180 y \$220/mes.

\section{Escenario de contingencia: OpenAI como principal, Mammouth como respaldo}
\label{sec:contingencia}

Si el objetivo real no es sustituir a OpenAI sino \textbf{tener con qué seguir
funcionando si un día OpenAI se cae}, el análisis cambia en lo económico pero no
en lo técnico:

\begin{itemize}
    \item El argumento de coste mensual (\$152 vs \$90) pierde fuerza, porque
    Mammouth no sería el motor de las 12.000 facturas/mes, solo estaría
    disponible para activarse en una caída puntual.
    \item El riesgo técnico (Structured Outputs estricto y visión) \textbf{sigue
    intacto}: si no funciona bien, se descubriría en el peor momento posible
    --- durante una caída real de OpenAI, con facturas acumulándose --- si no se
    ha probado antes. Debe verificarse en frío, no esperar a necesitarlo.
    \item El cambio de proveedor \textbf{no es automático hoy}: \texttt{build\_client()}
    lee las variables de entorno una sola vez al arrancar el servicio; pasar a
    Mammouth exige editar \texttt{.env} y reiniciar el servicio a mano. Si se
    quiere failover automático, hay que añadir lógica de detección de fallo y
    reintento que hoy no existe en el código.
\end{itemize}

\subsection{Alternativa a considerar: Azure OpenAI}

Para este objetivo concreto (resiliencia ante una caída del endpoint de
\texttt{openai.com}) existe una alternativa con \textbf{riesgo de compatibilidad
nulo}: \textbf{Azure OpenAI}. Son los mismos modelos (incluido \texttt{gpt-4.1}),
operados por Microsoft, con el mismo soporte garantizado de
\texttt{response\_format strict} y visión, y con SLA formal. El propio diseño de
\texttt{OPENAI\_BASE\_URL} en el proyecto parece anticipar precisamente este caso.
No protege frente a un problema sistémico del propio modelo, pero sí frente al
tipo de caída más habitual (infraestructura/API de OpenAI).

\section{¿Pagar más soluciona el problema de las cuotas?}

Con OpenAI directo el modelo es de \emph{pago por consumo puro}: mientras haya
saldo, cada petición se sirve al mismo modelo y a la misma tarifa, sin sorpresas.

Con Mammouth, según su propia documentación, el comportamiento \textbf{por
defecto} al agotar la cuota incluida no es ``cóbrame más y sigue igual'': es
\textbf{bajar de modelo automáticamente}. Mencionan la posibilidad de comprar
créditos adicionales de pago por consumo, pero:

\begin{itemize}
    \item No publican la tarifa de esos créditos adicionales.
    \item No confirman si comprarlos \emph{evita} la degradación de modelo o
    simplemente amplía la misma cuota, que puede volver a agotarse.
\end{itemize}

Antes de asumir que ``pagando se soluciona'', debería confirmarse por escrito con
su equipo comercial que el pago adicional impide el downgrade automático, y a qué
tarifa exacta --- especialmente relevante si el uso previsto es de contingencia
ante una caída real, donde no interesa descubrir esa letra pequeña en caliente.

\section{Cambios necesarios si se decide avanzar}

\begin{enumerate}
    \item \textbf{Prueba de compatibilidad real} antes de cualquier otra cosa:
    replicar la llamada exacta de \texttt{extract\_invoice\_with\_agent} (con
    \texttt{json\_schema strict}) y la de visión de \texttt{imagenes.py} contra el
    endpoint de Mammouth, con un caso de \emph{ground truth} conocido.
    \item \textbf{Código:} en principio basta cambiar \texttt{OPENAI\_BASE\_URL} y
    \texttt{OPENAI\_API\_KEY} en \texttt{.env}; no haría falta tocar
    \texttt{logic.py} ni \texttt{imagenes.py} salvo que fallen las pruebas del
    punto~1.
    \item \textbf{Detección de degradación de modelo:} como Mammouth no avisa al
    bajar de modelo, habría que comparar el campo \texttt{model} de la respuesta
    contra el esperado y alertar o frenar el pipeline si cambia, para no procesar
    facturas reales con un modelo peor sin saberlo.
    \item \textbf{Negociación comercial directa} con Mammouth: pedir tarifa de
    volumen publicada, confirmación por escrito de soporte de
    \texttt{response\_format strict} y visión, y SLA formal.
    \item \textbf{Comparar con alternativas} diseñadas desde el origen para
    consumo programático 24/7 con varios proveedores (Azure OpenAI, LiteLLM
    autoalojado, OpenRouter), en lugar de un producto de chat con API añadida
    encima.
\end{enumerate}

\section{Conclusión}

Como \textbf{sustituto} del proveedor principal, Mammouth.ai no aporta ventaja
clara sobre la integración directa con OpenAI que ya existe y está abstraída en
el código: introduce riesgos técnicos no confirmados sobre las dos capacidades
que el pipeline usa de forma crítica (Structured Outputs estricto y visión), un
modelo de cuotas pensado para chat interactivo que degrada silenciosamente la
calidad al alcanzar el límite, y un coste estimado un $\approx$69\% superior al
actual incluso en el escenario más favorable.

Como \textbf{respaldo de contingencia} ante una caída de OpenAI, el argumento
económico pierde peso, pero el riesgo técnico sigue siendo el mismo y debe
verificarse antes de necesitarlo. Para ese objetivo concreto, \textbf{Azure
OpenAI} es una alternativa con el mismo nivel de compatibilidad garantizada y
sin las incertidumbres de cuotas y soporte descritas en este informe.

\end{document}
